\section*{Hoofdstuk \ref{ch:b2:comparison} --- Vergelijking van functies}

\begin{solution}{exo:b2:comparison:1}
$\sqrt{x^2 + x + 1} = x\sqrt{1 + \tfrac1x + \tfrac{1}{x^2}} = x +
\frac12 + \frac38\cdot\frac1x + o\bigl(\frac1x\bigr)$ (binomiale
ontwikkeling: $\frac12 u - \frac18 u^2$ met $u = \frac1x +
\frac{1}{x^2}$ geeft $\frac{1}{2x} + \frac{1}{2x^2} -
\frac{1}{8x^2} = \frac{1}{2x} + \frac{3}{8x^2}$, en vermenigvuldig
daarna met $x$).

$\ln(x^2 + x) - 2\ln x = \ln\bigl(1 + \tfrac1x\bigr) = \frac1x -
\frac{1}{2x^2} + o\bigl(\frac{1}{x^2}\bigr)$.

Derde functie: ontwikkel elke factor,
\[
\frac{x + \sin x}{x - \ln x}
= \Bigl(1 + \frac{\sin x}{x}\Bigr)
\Bigl(1 + \frac{\ln x}{x} + \frac{(\ln x)^2}{x^2} +
O\Bigl(\frac{(\ln x)^3}{x^3}\Bigr)\Bigr).
\]
Orden de bijdragen op de schaal bij $+\infty$: $\frac{\ln x}{x}
\gg \frac{1}{x} \geq \bigl|\frac{\sin x}{x}\bigr| \gg \frac{(\ln
x)^2}{x^2}$. De twee termen na de hoofdterm $1$ zijn dus
$\frac{\ln x}{x}$, en daarna de begrensd oscillerende term
$\frac{\sin x}{x}$:
\[
\frac{x + \sin x}{x - \ln x}
= 1 + \frac{\ln x}{x} + \frac{\sin x}{x}
+ O\Bigl(\frac{(\ln x)^2}{x^2}\Bigr).
\]
\end{solution}

\begin{solution}{exo:b2:comparison:2}
Neem $f(t) = t^\alpha$ ($t \geq 1$).

$\alpha > -1$: divergentie, en volgens
\cref{thm:b2:comparison:seriesintegral}~(2) geldt $\sum_{k\leq n}
k^\alpha = \frac{n^{\alpha+1}}{\alpha+1} + C + o(1)$ als $\alpha <
0$ (daar daalt $f$); voor $\alpha \geq 0$ ($f$ stijgt) geeft
dezelfde insluiting met omgekeerde ongelijkheden $\sum_{k \leq n}
k^\alpha \sim \frac{n^{\alpha + 1}}{\alpha + 1}$.

$\alpha = -1$: $H_n = \ln n + \gamma + o(1)$
(\cref{ex:b2:comparison:harmonic}).

$\alpha < -1$: convergentie, met restsom $\sum_{k > n} k^\alpha
\sim \frac{n^{\alpha+1}}{-(\alpha+1)}$ volgens de insluiting (1)
(beide integraalgrenzen zijn met die waarde equivalent).
\end{solution}

\begin{solution}{exo:b2:comparison:3}
Zij $v_n = H_n - \ln n - \gamma \to 0$. Dan is
\[
v_n - v_{n+1} = \ln\frac{n+1}{n} - \frac{1}{n+1}
= \Bigl(\frac1n - \frac{1}{2n^2}\Bigr) - \Bigl(\frac1n -
\frac{1}{n^2}\Bigr) + O\Bigl(\frac{1}{n^3}\Bigr)
= \frac{1}{2n^2} + O\Bigl(\frac{1}{n^3}\Bigr),
\]
met $\frac{1}{n+1} = \frac1n - \frac{1}{n^2} + O(n^{-3})$. Omdat
$v_n \to 0$, geeft het telescoperen van de staart
\[
v_n = \sum_{k \geq n} (v_k - v_{k+1})
= \sum_{k\geq n} \Bigl(\frac{1}{2k^2} + O(k^{-3})\Bigr)
= \frac{1}{2n} + O\Bigl(\frac{1}{n^2}\Bigr),
\]
via \cref{thm:b2:comparison:seriesintegral}~(1) toegepast op
$t^{-2}$ (restsom $\sim \frac1n$, gehalveerd) en op $t^{-3}$.
Bijgevolg is $H_n = \ln n + \gamma + \frac{1}{2n} +
o(\frac1n)$.
\end{solution}

\begin{solution}{exo:b2:comparison:4}
Drie keer Stirling:
\[
\frac{(3n)!}{(n!)^3}
\sim \frac{\sqrt{6\pi n}\,(3n/\eu)^{3n}}
{(2\pi n)^{3/2}\,(n/\eu)^{3n}}
= \frac{\sqrt{6}\; 27^{\,n}}{2\pi n} \cdot
\frac{1}{\sqrt{2\pi n}}\cdot\sqrt{2\pi n}\;
= \frac{\sqrt3\,27^n}{2\pi n} .
\]
(Zorgvuldig: $\frac{\sqrt{6\pi n}}{(2\pi n)^{3/2}} =
\frac{\sqrt6}{(2\pi n)\sqrt{2\pi n}}\sqrt{\pi n} =
\frac{\sqrt3}{2\pi n}$.)

$\dfrac{n!}{n^n} \sim \sqrt{2\pi n}\,\eu^{-n}$.

$\sqrt[n]{n!} = \exp\bigl(\frac{\ln n!}{n}\bigr)$ met $\ln n! =
n\ln n - n + \frac12\ln(2\pi n) + o(1)$:
\[
\sqrt[n]{n!} = \exp\Bigl(\ln n - 1 + \frac{\ln(2\pi n)}{2n} +
o\Bigl(\frac{\ln n}{n}\Bigr)\Bigr)
= \frac{n}{\eu}\Bigl(1 + \frac{\ln(2\pi n)}{2n} +
o\Bigl(\frac{\ln n}{n}\Bigr)\Bigr).
\]
\end{solution}

\begin{solution}{exo:b2:comparison:5}
$g(x) = x^n + x - 1$ stijgt strikt op $\intcc{0}{1}$ van $-1$ naar
$1$: dus één nulpunt $x_n$. Omdat $x_n^n = 1 - x_n \in
\intoo{0}{1}$: zou $x_n \leq c < 1$ langs een deelrij, dan $x_n^n
\leq c^n \to 0$, dus $1 - x_n \to 0$: in strijd met $x_n \leq c$.
Bijgevolg is $x_n \to 1$.

Schrijf $x_n = 1 - \varepsilon_n$ met $\varepsilon_n \to 0^+$. De
vergelijking luidt $(1 - \varepsilon_n)^n = \varepsilon_n$, dat wil
zeggen
\[
n\ln(1 - \varepsilon_n) = \ln \varepsilon_n
\quad\Longrightarrow\quad
-n\varepsilon_n\bigl(1 + o(1)\bigr) = \ln\varepsilon_n .
\]
Dus $n\varepsilon_n = -\ln\varepsilon_n\,(1 + o(1)) \to +\infty$,
en opnieuw logaritmen nemen geeft $\ln n + \ln\varepsilon_n =
\ln(-\ln\varepsilon_n) + o(1)$. Omdat $\ln(-\ln \varepsilon_n) =
o(\ln(1/\varepsilon_n))$, volgt $\ln\varepsilon_n \sim -\ln n$,
waaruit $\varepsilon_n = \frac{-\ln\varepsilon_n}{n}(1 + o(1))
\sim \frac{\ln n}{n}$:
\[
x_n = 1 - \frac{\ln n}{n} + o\Bigl(\frac{\ln n}{n}\Bigr) .
\]
\end{solution}

\begin{solution}{exo:b2:comparison:6}
Exacte betrekking: $\cot y_n = x_n = n\pi + \frac\pi2 - y_n$ met
$y_n \sim \frac{1}{n\pi}$ (\cref{ex:b2:comparison:tan}). Ontwikkel
$\cot y = \frac1y - \frac y3 + O(y^3)$:
\[
\frac{1}{y_n} - \frac{y_n}{3} + O(y_n^3) = n\pi + \frac\pi2 - y_n
\quad\Longrightarrow\quad
\frac{1}{y_n} = n\pi + \frac\pi2 + O\Bigl(\frac1n\Bigr),
\]
(de termen $-y_n$ en $-\frac{y_n}{3}$ zijn $O(\frac1n)$). Inverteer:
\[
y_n = \frac{1}{n\pi}\cdot\frac{1}{1 + \frac{1}{2n} + O(n^{-2})}
= \frac{1}{n\pi}\Bigl(1 - \frac{1}{2n} + O\Bigl(\frac{1}{n^2}\Bigr)\Bigr)
= \frac{1}{n\pi} - \frac{1}{2n^2\pi} + O\Bigl(\frac{1}{n^3}\Bigr).
\]
Bijgevolg is
\[
x_n = n\pi + \frac{\pi}{2} - y_n
= n\pi + \frac\pi2 - \frac{1}{n\pi} + \frac{1}{2n^2\pi} +
o\Bigl(\frac{1}{n^2}\Bigr).
\]
\end{solution}

\begin{solution}{exo:b2:comparison:7}
Splits de twee grootste termen af:
\[
\sum_{k=0}^{n} k! = n! + (n-1)! + \sum_{k \leq n-2} k! ,
\qquad
\sum_{k\leq n-2} k! \leq (n-1)\,(n-2)! = (n-1)! .
\]
Dus $1 \leq \frac{1}{n!}\sum k! \leq 1 + \frac{2}{n}$: de limiet is
$1$. Verfijning: $\frac{(n-1)!}{n!} = \frac1n$, en de ruwe grens
$\sum_{k \leq n-2}k! \leq (n-1)!$ laat zich op dezelfde manier
aanscherpen: $\sum_{k\leq n-2} k! = (n-2)!\,(1 + O(\frac1n)) =
O\bigl(\frac{n!}{n^2}\bigr)$. Bijgevolg is
\[
\sum_{k=0}^{n} k! = n!\Bigl(1 + \frac1n + O\Bigl(\frac{1}{n^2}\Bigr)\Bigr).
\]
\end{solution}

\begin{solution}{exo:b2:comparison:8}
$(u_n)$ stijgt; was zij begrensd, dan zou zij naar een $\ell$
convergeren met $\ell = \ell + \frac1\ell$: absurd. Dus $u_n \to
\infty$.

Kwadraten: $u_{n+1}^2 = u_n^2 + 2 + u_n^{-2}$, dus
\[
u_n^2 = u_0^2 + 2n + \sum_{k=0}^{n-1} \frac{1}{u_k^2} .
\]
De som is $o(n)$ (de termen gaan naar $0$; Cesàro), dus $u_n^2 \sim
2n$ en $u_n \sim \sqrt{2n}$.

Verfijning: $\frac{1}{u_k^2} \sim \frac{1}{2k}$, dus levert de
vergelijking (\cref{thm:b2:comparison:seriesintegral}, of
equivalenten van partiële sommen van positieve reeksen)
$\sum_{k<n} u_k^{-2} \sim \frac12 \ln n$. Bijgevolg is
\[
u_n^2 = 2n + \frac{\ln n}{2}\,(1 + o(1)) + O(1)
\quad\Longrightarrow\quad
u_n = \sqrt{2n}\sqrt{1 + \frac{\ln n}{4n} + o\Bigl(\frac{\ln
n}{n}\Bigr)}
= \sqrt{2n}\Bigl(1 + \frac{\ln n}{8n} + o\Bigl(\frac{\ln
n}{n}\Bigr)\Bigr).
\]
\end{solution}

\begin{solution}{exo:b2:comparison:9}
Haal $n$ buiten haakjes en zet $t = \ln n$:
\[
S_n = \frac1n \sum_{k=1}^{n} \frac{1}{1 + t\,\frac kn} .
\]
Bij vaste $t$ is de som een riemannsom van $u \mapsto \frac{1}{1 +
tu}$ op $\intcc{0}{1}$; de functie is monotoon in $u$, dus wordt de
riemannsom ingesloten door de integraal, één maaswijdte verschoven:
\[
\int_0^1 \frac{\dd u}{1 + tu} - \frac1n
\leq S_n \leq \int_0^1 \frac{\dd u}{1 + tu} + \frac1n
\]
(de vergelijking van de riemannsommen van een monotone functie met
haar integraal, geldig voor elke $n$ met haar eigen $t = \ln n$).
Nu is $\int_0^1 \frac{\dd u}{1 + tu} = \frac{\ln(1 + t)}{t}$ en
$\frac1n = o\bigl(\frac{\ln t}{t}\bigr)$: bijgevolg
\[
S_n = \frac{\ln(1 + \ln n)}{\ln n} + O\Bigl(\frac 1n\Bigr)
\;\sim\; \frac{\ln\ln n}{\ln n} .
\]
\end{solution}

\begin{solution}{exo:b2:comparison:10}
Identiteit: $(\ln n)^{\ln n} = \eu^{\ln n\,\ln\ln n} =
\bigl(\eu^{\ln n}\bigr)^{\ln\ln n} = n^{\ln\ln n}$. Rangschikking:
vergelijk logaritmen. $\ln(n^2) = 2\ln n$; $\ln\bigl((\ln
n)^{\ln n}\bigr) = \ln n\ln\ln n$; $\ln(2^n) = n\ln2$; $\ln(n!) =
n\ln n - n + O(\ln n)$ (Stirling, of de ruwere insluiting $\ln n!
\sim n\ln n$); $\ln(n^n) = n\ln n$. Omdat $2\ln n = o(\ln n\ln\ln
n)$, $\ln n\ln\ln n = o(n)$, $n\ln 2 = o(n\ln n - n)$, en $n \ln n
- n \sim n\ln n$ terwijl $n! / n^n \to 0$ (het verschil van de
logaritmen is $-n + O(\ln n) \to -\infty$), volgt
\[
n^2 = o\bigl((\ln n)^{\ln n}\bigr),\quad
(\ln n)^{\ln n} = o(2^n),\quad
2^n = o(n!),\quad
n! = o(n^n).
\]
(Bij elke stap: het verschil van de logaritmen gaat naar
$+\infty$, dus gaat de verhouding naar $0$.)
\end{solution}

\begin{solution}{exo:b2:comparison:11}
Splits $\frac1{k^2} = \frac1{k(k+1)} + \frac1{k^2(k+1)}$ en sommeer
voor $k > n$:
\[
\sum_{k>n}\frac1{k^2} = \frac1{n+1} +
\sum_{k>n}\frac{1}{k^2(k+1)} ,
\]
waarbij de eerste som exact telescopeert ($\frac1{k(k+1)} = \frac1k
- \frac1{k+1}$). Voor de tweede: $\frac{1}{k^2(k+1)} = \frac1{k^3}
+ O\bigl(\frac1{k^4}\bigr)$ (want $\frac{1}{k^2(k+1)} -
\frac1{k^3} = \frac{-1}{k^3(k+1)}$), en de vergelijking met de
integraal geeft $\sum_{k>n}\frac1{k^3} = \frac1{2n^2} +
O\bigl(\frac1{n^3}\bigr)$ en $\sum_{k>n}\frac1{k^4} =
O\bigl(\frac1{n^3}\bigr)$. Bijgevolg is
\[
\sum_{k>n}\frac1{k^2}
= \frac1{n+1} + \frac{1}{2n^2} + O\Bigl(\frac1{n^3}\Bigr)
= \frac1n - \frac1{n^2} + \frac{1}{2n^2} +
O\Bigl(\frac1{n^3}\Bigr)
= \frac1n - \frac{1}{2n^2} + O\Bigl(\frac1{n^3}\Bigr),
\]
met $\frac1{n+1} = \frac1n - \frac1{n^2} +
O\bigl(\frac1{n^3}\bigr)$.
\end{solution}

\begin{solution}{exo:b2:comparison:12}
$(u_n)$ stijgt; was zij begrensd, dan zou zij naar een eindige
$\ell$ convergeren met $\ell = \ell + \eu^{-\ell}$: onmogelijk. Dus
$u_n \to \infty$. Zij $v_n = \eu^{u_n} \to \infty$; dan is
\[
v_{n+1} = \eu^{u_n + \eu^{-u_n}} = v_n\,\eu^{1/v_n}
= v_n\Bigl(1 + \frac1{v_n} + \frac1{2v_n^2} +
O\bigl(v_n^{-3}\bigr)\Bigr)
= v_n + 1 + \frac{1}{2v_n} + O\bigl(v_n^{-2}\bigr).
\]
Het sommeren van $v_{k+1} - v_k = 1 + O(1)$ geeft eerst $v_n = n +
O(n)$, dus $v_n \geq cn$ vanaf zekere rang; opnieuw sommeren met
$\frac1{2v_k} = O(\frac1k)$ geeft $v_n = n + O(\ln n)$. Nog een
ronde: $\frac{1}{2v_k} = \frac{1}{2k}\bigl(1 +
O\bigl(\tfrac{\ln k}k\bigr)\bigr)$, dus
\[
v_n = n + \sum_{k<n}\frac1{2k} + O(1) = n + \frac{\ln n}2 +
O(1).
\]
Ten slotte $u_n = \ln v_n = \ln n + \ln\Bigl(1 + \frac{\ln n}{2n} +
O\bigl(\tfrac1n\bigr)\Bigr) = \ln n + \frac{\ln n}{2n} +
O\bigl(\tfrac1n\bigr)$.
\end{solution}

\begin{solution}{pb:b2:comparison:1}
\textbf{1.} Aftrekken van de twee ontwikkelingen geeft $\sum_i (c_i
- c_i')\varphi_i = o(\varphi_k)$. Verschilt een coëfficiënt, zij
$i_0$ de eerste: deling door $\varphi_{i_0}$ en het gebruik van
$\varphi_j = o(\varphi_{i_0})$ voor $j > i_0$ geeft $c_{i_0} -
c_{i_0}' = o(1)$, dus nul: tegenspraak. Voor de ontwikkeling: met
$u = \frac{\ln x}x \to 0$ is
\[
\frac{1}{x - \ln x} = \frac1x\cdot\frac{1}{1 - u}
= \frac1x\bigl(1 + u + u^2 + O(u^3)\bigr)
= \frac1x + \frac{\ln x}{x^2} + \frac{(\ln x)^2}{x^3} +
o\Bigl(\frac{(\ln x)^2}{x^3}\Bigr).
\]
Er verschijnt geen term $\frac c{x^2}$ omdat de ontwikkeling een
meetkundige reeks in $u = \frac{\ln x}{x}$ is: elke term draagt
voorbij de eerste evenveel machten van $\ln x$ als van $\frac1x$;
de sport $\frac1{x^2}$ van de schaal (de coëfficiënt bij $(\ln
x)^0$) ontbreekt eenvoudigweg, met coëfficiënt $0$.

\textbf{2.} $f(x) = \eu^x + x$ is \omterm{def:b2:metric:continuity}{continu} en strikt stijgend, met
limieten $-\infty$ en $+\infty$: dus een bijectie $\R \to \R$,
zodat $x_n = f^{-1}(n)$ bestaat en eenduidig is, en $x_n \to
+\infty$ ($f^{-1}$ stijgt naar $+\infty$). Uit $\eu^{x_n} = n -
x_n$ volgt $x_n = \ln(n - x_n) \leq \ln n$, dus $x_n/n \to 0$ en
$x_n = \ln n + \ln(1 - x_n/n) = \ln n + o(1) \sim \ln n$.

\textbf{3.} Schrijf $u_n = x_n/n$. Tweede ronde: $u_n = \frac{\ln n
+ o(1)}{n}$, dus
\[
x_n = \ln n + \ln(1 - u_n) = \ln n - u_n + O(u_n^2)
= \ln n - \frac{\ln n}{n} + o\Bigl(\frac{\ln n}n\Bigr).
\]
Derde ronde: nu is $u_n = \frac{\ln n}{n} - \frac{\ln n}{n^2} +
o\bigl(\frac{\ln n}{n^2}\bigr)$, en $\ln(1 - u_n) = -u_n -
\frac{u_n^2}2 + O(u_n^3)$:
\[
x_n = \ln n - \frac{\ln n}n + \frac{\ln n}{n^2}
- \frac{(\ln n)^2}{2n^2} + o\Bigl(\frac{(\ln
n)^2}{n^2}\Bigr)
= \ln n - \frac{\ln n}{n} - \frac{(\ln n)^2}{2n^2} +
o\Bigl(\frac{(\ln n)^2}{n^2}\Bigr),
\]
waarbij de term $\frac{\ln n}{n^2}$ in $o\bigl(\frac{(\ln
n)^2}{n^2}\bigr)$ wordt opgeslokt.

\textbf{4.} Bij $n = 1000$: $\ln 1000 \approx 6.90776$ (fout
$7\cdot10^{-3}$); met twee termen $6.90085$ (fout
$2\cdot10^{-5}$); met drie termen $6.90082$ (fout onder
$10^{-5}$), tegenover $x_{1000} \approx 6.90083$. Elke ronde koopt
ruwweg de voorspelde factor $\frac{\ln n}{n}$.

\textbf{5.} Twee partiële integraties, vanaf rechts: met
$\frac{\dd}{\dd t}\bigl[\tfrac12t(1-t)\bigr] = \tfrac12 - t$ en
$t(1-t)$ dat in beide eindpunten verdwijnt, is
\[
\frac12\int_0^1 t(1-t)g''(t)\dd t
= -\int_0^1\Bigl(\frac12 - t\Bigr)g'(t)\dd t
= -\Bigl[\Bigl(\frac12 - t\Bigr)g\Bigr]_0^1 - \int_0^1 g
= \frac{g(0) + g(1)}2 - \int_0^1 g .
\]
Herschikt is dat de gevraagde identiteit.

\textbf{6.} Bereken de toename en pas dan vraag 5 toe op $g(t) =
f(n + t)$:
\begin{align*}
E_{n+1} - E_n
&= f(n{+}1) - \int_n^{n+1}\!f - \frac{f(n{+}1) - f(n)}2 \\
&= \frac{f(n) + f(n{+}1)}2 - \int_n^{n+1}\!f
= \frac12\int_0^1 t(1-t)f''(n+t)\dd t .
\end{align*}
Omdat $0 \leq t(1-t) \leq \frac14$, is $\abs{E_{n+1} - E_n} \leq
\frac18\int_n^{n+1}\abs{f''}$, en de som daarvan over $n$
convergeert wegens de hypothese: $(E_n)$ convergeert (absoluut
sommeerbare toenamen) naar zekere $E$, met
\[
\abs{E - E_n} \leq \sum_{k\geq n}\abs{E_{k+1} - E_k} \leq
\frac18\int_n^\infty\abs{f''} .
\]

\textbf{7.} Voor $f(t) = \frac1t$ is $f''(t) = \frac2{t^3}$ en
$\int_1^\infty\abs{f''} = 1 < \infty$. Vraag 6 geeft
\[
H_n = \ln n + \frac{1 + \frac1n}{2} + E + (E_n - E)
= \ln n + \Bigl(E + \frac12\Bigr) + \frac1{2n} +
\varepsilon_n,
\]
met $\abs{\varepsilon_n} = \abs{E_n - E} \leq
\frac18\int_n^\infty\frac{2\dd t}{t^3} = \frac1{8n^2}$.
Vergelijking met $H_n = \ln n + \gamma + o(1)$
(\cref{ex:b2:comparison:harmonic}) identificeert $E + \frac12 =
\gamma$.

\textbf{8.} Uit de formule voor de toename in vraag 6:
\[
\varepsilon_n = E_n - E = -\sum_{k\geq n}\frac12\int_0^1
t(1-t)\,\frac{2\,\dd t}{(k+t)^3}
= -\sum_{k \geq n}\Bigl(\frac1{k^3}\int_0^1t(1-t)\dd t +
O\Bigl(\frac1{k^4}\Bigr)\Bigr),
\]
met $\frac{1}{(k+t)^3} = \frac1{k^3} + O\bigl(\frac1{k^4}\bigr)$
uniform voor $t \in \intcc01$. Met $\int_0^1 t(1-t) = \frac16$ en
$\sum_{k\geq n}\frac1{k^3} \sim \frac{1}{2n^2}$
(\cref{thm:b2:comparison:seriesintegral}):
\[
\varepsilon_n = -\frac16\cdot\frac{1}{2n^2} +
o\Bigl(\frac1{n^2}\Bigr) = -\frac{1}{12n^2} +
o\Bigl(\frac{1}{n^2}\Bigr).
\]

\textbf{9.} Voor $f = \ln$ is $f''(t) = -\frac1{t^2}$, absoluut
integreerbaar. Vraag 6 geeft
\[
\ln n! = \int_1^n\ln t\,\dd t + \frac{\ln n}2 + E + O\Bigl(
\frac1{8}\int_n^\infty\frac{\dd t}{t^2}\Bigr)
= \Bigl(n + \frac12\Bigr)\ln n - n + 1 + E +
O\Bigl(\frac1n\Bigr),
\]
dus $d_n = 1 + E + O\bigl(\frac1n\bigr)$: de convergentie van
$(d_n)$ --- stap 1 van \cref{thm:b2:comparison:stirling} --- plus
de snelheid $O(1/n)$. (De waarde van de limiet volgens Stirling
geeft $E = \ln\sqrt{2\pi} - 1$.)

\textbf{10.} Voor $f(t) = t^{-1/2}$ is $f''(t) = \frac34
t^{-5/2}$, absoluut integreerbaar. Vraag 6 geeft
\[
\sum_{k=1}^n \frac1{\sqrt k}
= 2\sqrt n - 2 + \frac{1 + \frac1{\sqrt n}}2 + E +
O\bigl(n^{-3/2}\bigr)
= 2\sqrt n + c + \frac{1}{2\sqrt n} +
O\bigl(n^{-3/2}\bigr),
\]
met $c = E - \frac32$. Bij $n = 10^4$: $2\sqrt n = 200$, $c
\approx -1.46035$, $\frac1{2\sqrt n} = 0.005$: voorspeld
$198.54465$, en inderdaad is $\sum_{k\leq10^4}k^{-1/2} =
198.544645\dots$ --- drie termen, zeven cijfers.

\textbf{11.} $t \mapsto t\ln t$ is \omterm{def:b2:metric:continuity}{continu} en strikt stijgend op
$\intco1\infty$ (afgeleide $\ln t + 1 \geq 1$), van $0$ tot
$+\infty$: dus bestaat er een unieke $x_n$, en $x_n \to \infty$
(anders zou $x_n\ln x_n$ begrensd blijven). Logaritmen nemen in
$x_n\ln x_n = n$ geeft $\ln x_n + \ln\ln x_n = \ln n$; omdat
$\ln\ln x_n = o(\ln x_n)$, geeft deling door $\ln x_n$ dat
$\frac{\ln n}{\ln x_n} \to 1$: dus $\ln x_n \sim \ln n$.

\textbf{12.} Uit $x_n = \frac{n}{\ln x_n}$ en $\ln x_n \sim \ln n$
volgt $x_n \sim \frac{n}{\ln n}$. Volgende ronde: $\ln\ln x_n =
\ln\bigl(\ln n\,(1 + o(1))\bigr) = \ln\ln n + o(1)$, dus $\ln x_n =
\ln n - \ln\ln n + o(1)$ en
\[
x_n = \frac{n}{\ln n - \ln\ln n + o(1)}
= \frac{n}{\ln n}\cdot\frac{1}{1 - \frac{\ln\ln n +
o(1)}{\ln n}}
= \frac{n}{\ln n}\Bigl(1 + \frac{\ln\ln n}{\ln n} +
o\Bigl(\frac{\ln\ln n}{\ln n}\Bigr)\Bigr).
\]

\textbf{13.} Bij $n = 10^6$: $\frac{n}{\ln n} \approx 72\,382$
($18\%$ ernaast), met twee termen $\approx 86\,140$ ($1.9\%$
ernaast), tegenover de werkelijke $x \approx 87\,848$. De winst per
ronde is slechts de factor $\frac{\ln\ln n}{\ln n} \approx
\frac{2.63}{13.8} \approx 0.19$: logaritmische schalen convergeren
tergend langzaam --- een feit van het leven overal waar
priemgetallen in het spel zijn.

\textbf{14.} Er zijn precies $n$ priemgetallen $\leq p_n$ (namelijk
$p_1, \dots, p_n$): dus $\pi(p_n) = n$. De priemgetalstelling
(aangenomen; volume van bachelorjaar 3) geeft $n = \pi(p_n) \sim
\frac{p_n}{\ln p_n}$, dat wil zeggen $p_n \sim n\ln p_n$: dat is de
vergelijking $x\ln x \approx n$, achterstevoren gelezen. Logaritmen
nemen: $\ln p_n = \ln n + \ln\ln p_n + o(1)$, en $\ln\ln p_n =
o(\ln p_n)$ dwingt $\ln p_n \sim \ln n$ af, net als in vraag 11.
Terugsubstitueren geeft
\[
p_n \sim n\ln p_n = n\,\ln n\,\frac{\ln p_n}{\ln n} \sim n\ln
n .
\]

\textbf{15.} (a) Leg $\varepsilon > 0$ vast; voor grote $k$ is $(1
- \varepsilon)k\ln k \leq p_k \leq (1 + \varepsilon)k\ln k$. Met
de vergelijking met de stijgende $t\ln t$ (insluiting van het type
\cref{thm:b2:comparison:seriesintegral}) is $\sum_{k\leq n}k\ln k =
\int_1^n t\ln t\,\dd t + O(n\ln n) = \frac{n^2\ln n}2 -
\frac{n^2}4 + O(n\ln n) \sim \frac{n^2\ln n}2$. Bijgevolg is
$\sum_{k\leq n}p_k = \frac{n^2\ln n}{2}(1 + O(\varepsilon) +
o(1))$ voor elke $\varepsilon$: dus $\sum_{k\leq n}p_k \sim
\frac{n^2\ln n}2$. (b) Volgens de priemgetalstelling is onder de
gehele getallen tot $10^{100}$ een fractie $\sim \frac{1}{\ln
10^{100}} = \frac1{230.26\dots}$ priem: een uniform gekozen geheel
getal van $100$ cijfers is priem met kans ongeveer
$\frac1{230}$.

\textbf{16.} Op $\intoo{n\pi}{n\pi + \frac\pi2}$ is $g(x) = \tan x
- \frac1x$ \omterm{def:b2:metric:continuity}{continu} en strikt stijgend ($g' = 1 + \tan^2x +
\frac1{x^2} > 0$), met $g \to -\frac1{n\pi} < 0$ aan de
linkerkant en $g \to +\infty$ aan de rechterkant: dus precies één
nulpunt $x_n$. Omdat $\tan z_n = \tan x_n = \frac1{x_n} \to 0$ met
$z_n \in \intoo{0}{\frac\pi2}$, volgt $z_n \to 0^+$.

\textbf{17.} Er geldt $\tan z_n \sim z_n$ en $\frac1{x_n} \sim
\frac1{n\pi}$: dus $z_n \sim \frac1{n\pi}$.

\textbf{18.} Er is $z_n = \arctan\frac1{x_n}$ met $\arctan u = u +
O(u^3)$. Met $z_n = O(\frac1n)$:
\[
\frac1{x_n} = \frac{1}{n\pi}\cdot\frac1{1 + \frac{z_n}{n\pi}}
= \frac1{n\pi} - \frac{z_n}{n^2\pi^2} +
O\Bigl(\frac1{n^4}\Bigr)
= \frac1{n\pi} + O\Bigl(\frac1{n^3}\Bigr),
\]
dus $z_n = \frac1{n\pi} + O\bigl(\frac1{n^3}\bigr)$: de sport
$\frac{c}{n^2}$ heeft coëfficiënt $0$, omdat de eerste correctie op
$\frac1{x_n}$ zelf van de grootte $\frac{z_n}{n^2} = O(n^{-3})$
is.

\textbf{19.} Vul $z_n = \frac1{n\pi} + O(n^{-3})$ in de vorige
formule in:
\[
\frac{1}{x_n} = \frac{1}{n\pi} - \frac{1}{n^3\pi^3} +
O\Bigl(\frac1{n^5}\Bigr),
\]
en dan $z_n = \arctan\frac1{x_n} = \frac1{x_n} -
\frac{1}{3}\Bigl(\frac1{x_n}\Bigr)^3 + O\Bigl(\frac1{n^5}\Bigr)
= \frac1{n\pi} - \frac{1}{n^3\pi^3} - \frac{1}{3n^3\pi^3} +
O\Bigl(\frac1{n^5}\Bigr)$:
\[
x_n = n\pi + \frac{1}{n\pi} - \frac{4}{3\pi^3n^3} +
O\Bigl(\frac1{n^5}\Bigr).
\]

\textbf{20.} Bij $n = 3$: met één term $9.53088$, met drie termen
$9.52929$, tegenover het werkelijke nulpunt $9.52933$: fouten
$1.5\cdot10^{-3}$ en $5\cdot10^{-5}$. Contrast: bij $\tan x = x$
moet het nulpunt $\tan$ enorm maken, dus kruipt het tegen het
\emph{rechter} uiteinde $n\pi + \frac\pi2$ van het venster aan, op
afstand $\sim\frac1{n\pi}$ vóór de asymptoot; bij $x\tan x = 1$
moet het nulpunt $\tan$ minuscuul maken, dus ligt het net voorbij
het \emph{linker} uiteinde $n\pi$, op afstand $\sim\frac1{n\pi}$ na
het nulpunt. Dezelfde methode, gespiegelde geografie.

\textbf{21.} Er geldt $\sin u < u$ op $\intoo0\pi$, en $\sin$
beeldt $\intoo0\pi$ af in $\intoc01 \subseteq \intoo0\pi$: na één
stap ligt $u_1 \in \intoc{0}{1}$, en daarna daalt $(u_n)$ en is zij
van onderen door $0$ begrensd: zij convergeert dus naar een vast
punt van $\sin$, dat wil zeggen naar $0$. Ontwikkeling: $\sin u =
u(1 - \frac{u^2}6 + o(u^2))$, dus
\[
\frac{1}{u_{n+1}^2} - \frac1{u_n^2}
= \frac{1}{u_n^2}\Bigl(\bigl(1 - \tfrac{u_n^2}6 +
o(u_n^2)\bigr)^{-2} - 1\Bigr)
= \frac{1}{u_n^2}\Bigl(\frac{u_n^2}{3} + o(u_n^2)\Bigr)
\longrightarrow \frac13 .
\]

\textbf{22.} Volgens Cesàro (volume van bachelorjaar 1)
convergeert het gemiddelde van de toenamen naar dezelfde limiet:
\[
\frac{1}{n}\cdot\frac{1}{u_n^2}
= \frac1n\Bigl(\frac1{u_0^2} + \sum_{k=0}^{n-1}
\Bigl(\frac1{u_{k+1}^2} - \frac1{u_k^2}\Bigr)\Bigr)
\longrightarrow \frac13 ,
\]
dus $u_n^2 \sim \frac3n$ en, omdat alle termen positief zijn, $u_n
\sim \sqrt{3/n}$.

\textbf{23.} Volgens vraag 7 is $\abs{H_n - \ln n - \gamma -
\frac1{2n}} \leq \frac1{8n^2}$. Bij $n = 10^6$ is die grens
$\frac{1}{8\cdot10^{12}} = 1.25\cdot10^{-13}$: drie berekende
termen leveren de harmonische som van een miljoen termen tot op
dertien cijfers, met een \omterm{def:b2:metric:complete}{volledig} rigoureus foutcertificaat --- en
daar draait een asymptotische formule met expliciete restterm
precies om.

\textbf{24.} (a) Waar: $\ln u_n - \ln v_n = \ln\frac{u_n}{v_n} \to
0$ terwijl $\ln v_n \to +\infty$, dus gaat de verhouding van de
logaritmen naar $1$. (b) Onwaar: $u_n = n + 1 \sim v_n = n$, maar
$\eu^{u_n}/\eu^{v_n} = \eu \neq 1$. Equivalentie verdraagt
additieve fouten $o(1)$ in de exponent, geen $O(1)$. (c) Onwaar:
$f(x) = x + \sin(x^2) \sim g(x) = x$ bij $+\infty$, maar $f'(x) = 1
+ 2x\cos(x^2)$ oscilleert onbegrensd terwijl $g' = 1$: de
afgeleiden van equivalente functies hoeven in het geheel niet
vergelijkbaar te zijn.

\textbf{25.} De lus van \cref{met:b2:comparison:implicit} liep drie
keer identiek: lokaliseer het nulpunt, haal een ruwe term eruit,
voer die terug voor de volgende orde --- op $\eu^x + x = n$ (Deel
I), op $x\ln x = n$ (Deel III) en op $x\tan x = 1$ (Deel IV). De
trapeziumcorrectie tilt de vergelijking van reeks en integraal op
van ``het verschil convergeert'' naar een expliciete term
$\frac{f(1) + f(n)}2$ met een gecertificeerde restterm
$O(\int_n^\infty \abs{f''})$ --- constanten en foutbalken in plaats
van louter convergentie. De brug naar de priemgetallen is zuivere
omkering: de priemgetalstelling zegt $\pi(x)\ln x \approx x$, dus
lost $p_n$, gedefinieerd door $\pi(p_n) = n$, een vergelijking
$x\ln x = n$ op --- en erft haar asymptotiek. Regel (a) van vraag
24 legitimeerde elke overgang van $u_n \sim v_n$ naar $\ln u_n \sim
\ln v_n$ (vragen 11 en 14); dat (b) onwaar is, is de reden dat we
equivalenties nooit exponentiëren. De toppen: de formule van
Euler--Maclaurin tot op de eerste orde (vraag 6), en de
asymptotische wet $p_n \sim n\ln n$ voor het $n$-de priemgetal
(vraag 14).
\end{solution}
